\documentclass[11pt]{article}
\usepackage{amsmath,amssymb,amsthm,booktabs}
\usepackage[margin=1in]{geometry}

\newtheorem{definition}{Definition}
\newtheorem{theorem}{Theorem}

\title{Step 180: Carrier-Typed Matrix-Element Terminality}
\author{Six Birds Foundations III}
\date{2026-05-15}

\begin{document}
\maketitle

\section{Verdict}

\[
\boxed{\texttt{final\_verdict = V\_CTMT\_formalized}.}
\]

This is case-3 synthesis: a foundational typed condition is recognized from a
track instance and formalized for reuse across carriers and future tracks.

\section{Definition}

\begin{definition}[Carrier-Typed Matrix-Element Terminality]
A residual closure attempt on a typed carrier \(C\) exhibits
\emph{Carrier-Typed Matrix-Element Terminality} (CTMT) when, after applying the
available operational identities and structural reductions, the closure
question reduces to deciding a carrier-native matrix-element family
\[
M(\alpha,\beta,\gamma)
=
\langle a_\alpha,P_\beta b_\gamma\rangle,
\]
where \(a_\alpha\) and \(b_\gamma\) are carrier-native vectors, \(P_\beta\) is
a carrier-native projection, transport, quotient, or operational kernel, and
the inherited records do not supply enough kernel/transport data to compute
\(M\).  The \(M\)-decision must not already be target-equivalent to the desired
conclusion, and public-shadow or finite-window data may not be promoted without
a typed bridge theorem.
\end{definition}

\section{Resolution Modes}

\[
\begin{array}{l|p{0.68\linewidth}}
\toprule
\text{Mode} & \text{Meaning}\\
\midrule
\text{CTMT-stuck} &
The \(M\)-family is well typed but a kernel or transport theorem is missing.
The \(M\)-family becomes the exact external-content interface.\\
\text{CTMT-foreclosed-numerical} &
A different lawful operational chain computes at least one \(M\)-value
nonzero with a sharp lower bound.  The route is foreclosed under inherited
normalization.\\
\text{CTMT-bridge-failure} &
The attempted \(M\)-family crosses carriers, but the carrier comparison data
are missing, so the carrier-pivot fails under inherited records.\\
\bottomrule
\end{array}
\]

\section{Four Instances}

\[
\begin{array}{c|p{0.17\linewidth}|p{0.24\linewidth}|p{0.18\linewidth}|p{0.18\linewidth}}
\toprule
\text{Instance} & \text{Carrier} & M\text{-form} & \text{Projection / bridge} & \text{Verdict}\\
\midrule
A &
Burnol/Sonine &
\(\langle\eta_i,C_\ell\eta_j\rangle\) &
\(P_\eta\), \(P_\infty\) via \(K_\infty^{op}\) &
\texttt{V\_branch\_a\_stuck\_at\_kappa}\\
B &
Burnol/Sonine &
\(c_{ij}(\ell)=\langle e_i,C_\ell e_j\rangle\) &
\(P_\infty\) via \(K_\infty^{op}\), pulled evaluators via \(\kappa\) &
\texttt{V\_kappa\_classical\_theorem\_needed}\\
C &
Burnol/Sonine &
\(L_{\rho,k}(G)=\langle M_\zeta G,P_\infty y_{\rho,k}\rangle\) &
\(P_\infty\) via \(K_\infty^{op}\) &
\texttt{V\_branch\_c\_foreclosed\_generic}\\
H6 &
Hecke to Burnol/Sonine &
attempted cross-carrier pairing &
carrier comparison / descent bridge &
\texttt{V-NC}\\
\bottomrule
\end{array}
\]

Instances A and B are CTMT-stuck at the same Burnol projected-kernel theorem
for \(\kappa\).  Instance C is CTMT-foreclosed-numerical: four nonzero
\(L\)-values with \(|L|\ge0.11\) foreclose the full-carrier Branch C route.
Instance H6 is CTMT-bridge-failure: the scalar identity \(L_{\mathbb Q}(s,1)
=\zeta(s)\) does not identify the carriers, probes, energy, or ledgers.

\section{Consequence Theorem}

\begin{theorem}[CTMT Consequence]
When a typed carrier exhibits CTMT, the framework's typed contribution is the
carrier-native matrix-element statement \(M\), including its carrier, vectors,
projection or transport, and resolution mode.  The outcome is exactly one of
the following typed statuses:
\[
\text{CTMT-stuck},\qquad
\text{CTMT-foreclosed-numerical},\qquad
\text{CTMT-bridge-failure}.
\]
Consequently, the external-content interface produced by the framework is the
\(M\)-statement, not a vague request for external operator theory.
\end{theorem}

\begin{proof}
The definition requires that all available typed reductions have already been
applied.  If the \(M\)-family remains well-defined but uncomputed, the missing
kernel or transport theorem is precisely the unresolved external content.  If
an alternate lawful operational chain computes a nonzero \(M\)-value, the
route's asserted vanishing or factorization fails under the inherited
normalization.  If the carriers cannot be compared, then no cross-carrier
matrix element is defined, so the descent or pivot fails.  These alternatives
exhaust the documented Step 175--179 instances and are carrier-typed rather
than public-shadow claims.
\end{proof}

\section{Corpus Candidate}

CTMT is flagged for future foundational-corpus inclusion.  Suggested target:
\[
\texttt{anti\_loc/CTMT.md}.
\]
This step writes only the standalone candidate
\[
\texttt{CTMT\_typed\_condition.md}
\]
inside the Step 180 artifact directory.

\section{Nonclaim Boundary}

CTMT does not prove RH, does not close \(\Xi_{\mathrm{BC}}\), does not close
any branch by itself, and does not weaken retained no-gos.  It is a typed
diagnostic and external-interface condition.

\end{document}
